% =========================================================================
% Chapter 6: Discussion
% =========================================================================

\chapter{Discussion}
\label{ch:discussion}

\section{Extract and phytochemistry}
\label{sec:disc_yield}

The recorded aqueous yield (24.13\% w/w) is larger than the dichloromethane and methanol yields of \citet{drakhshaan2026}. Water extracts glycosides, sugars, and mucilage together with phenolics \citep{cowan1999}. A larger crude mass does not imply a higher titre of the lipophilic terpenoids that dominated that dichloromethane GC-MS profile. Qualitative tests were positive for phenolics, flavonoids, alkaloids, saponins, tannins, terpenoids, phytosterols, quinones, and carbohydrates. Those classes match genus-level reviews \citep{zhao2011,skala2023} and polar HPLC markers reported from a methanol extract of the same species \citep{drakhshaan2026}. None of the associated mechanisms was measured here.

\section{What the zones do and do not show}
\label{sec:disc_antibacterial}

Distilled-water controls were zero on every plate, so clearing around extract discs is not a solvent artefact. Gentamicin 5 discs produced the largest zones in each species, as expected for a purified aminoglycoside on a viable lawn \citep{bauer1966,clsi2024}. Extract zones were smaller and species-dependent. Diameters were read inclusive of the discs.

\textit{E.\ coli} and \textit{K.\ pneumoniae} both showed larger means at C2 (\qty{100}{\milli\gram\per\milli\litre}) than at C1 (\qty{50}{\milli\gram\per\milli\litre}), consistent with more diffusible mass at the higher load. \textit{P.\ aeruginosa} did not: C1 exceeded C2. Possible explanations, none tested here, include incomplete wetting of the higher load, plate-to-plate variation in lawn density, or a measurement error. The last cannot be checked because no \textit{P.\ aeruginosa} inhibition photograph exists. The inverse pair is therefore reported as a non-monotonic observation, not as evidence that a lower dose is more potent.

Zone diameter on LB agar is a screening endpoint \citep{balouiri2016,valgas2007}. It is not an MIC and not a CLSI interpretive category \citep{andrews2001,vanvuuren2019}. Direct numerical comparison with \citet{drakhshaan2026} is limited by medium, inoculum, extract chemistry, and the fact that that study used organic solvents and well diffusion.

\section{\textit{Klebsiella pneumoniae} and indexed literature}
\label{sec:disc_klebsiella}

\citet{drakhshaan2026} tested \textit{E.\ coli} and \textit{P.\ aeruginosa} but not \textit{K.\ pneumoniae}. \citet{ji2026} used Kashmir \textit{D.\ inermis} as a reducing agent for ZnO nanoparticles assayed against \textit{S.\ aureus} and \textit{E.\ coli}. Kashmir disc-diffusion work on other temperate plants has included the same three Gram-negative genera \citep{guna2018,nabi2022}, but not this aqueous teasel extract. The present zones against \textit{K.\ pneumoniae} therefore fill an indexed-literature gap for this species--extract combination. They do not identify a carbapenemase genotype or a drug lead.

\section{Method limitations}
\label{sec:disc_limits}

\begin{enumerate}
  \item LB agar is not Mueller--Hinton agar \citep{clsi2024,balouiri2016}.
  \item Stocks have no ATCC numbers.
  \item The positive control is a GEN~5 disc (5\,\textmu g), not a 10\,\textmu g CLSI quality-control claim.
  \item No \textit{P.\ aeruginosa} inhibition plate was filed.
  \item No HPLC, GC-MS, MIC, or cytotoxicity assay was done on this aqueous extract.
\end{enumerate}

These limits keep the claims to a polar extract, qualitative phytochemistry, and triplicate disc-diffusion zones on LB agar.
